%% ===================================================================
\section{\vrev{주요 연구 성과}}
\label{sec:conclusion-achievements}
%% ===================================================================

\revdel{본 연구의 5대 핵심 성과를 다음과 같이 정리한다.}\,\revadd{본 항은 5장 전체에서 도출된 5대 핵심 성과를 압축 요약 형태로 제시한다. 각 성과의 상세 논의는 학술적 측면은 \secref{sec:academic-contrib}, 실무적 측면은 \secref{sec:practical-contrib}을 참조하라.}


%% -------------------------------------------------------------------
\chg{첫째 성과는 물리적 경계 보장을 갖춘 개별 시그모이드 바운딩 $\DRTO$ 모델이다.}
\label{subsec:achiev-model}
\chg{개별 시그모이드 바운딩을 통해 $\DRTO \in (0, R_{\max})$를 보장하고, $C^{\infty}$ 연속성과 자연스러운 교호작용 포착을 함께 달성하였다(상세는 \subsecref{subsec:academic-theoretical}).}

\fdel{업무연속성관리 분야 최초로, 관측성과 불확실성을 단일 수리 모델에서
통합하는 개별 시그모이드 바운딩 $\DRTO$ 모델을 정립하였다.
수정 시그모이드 함수 $S(z) = 2/(1 + e^{-z}) - 1$을 기저 함수로 채택하여,
관측성 효과($E_{O,\text{bnd}}$)와 불확실성 효과($E_{U,\text{bnd}}$)를
각각 고유한 물리적 범위($(-R_0,\; 0]$과 $[0,\; R_{\max} - R_0)$)에서
독립적으로 바운딩하였다. 이 모델은 물리적 경계 보장($\DRTO \in (0, R_{\max})$),
무한 미분 가능성($C^{\infty}$), 교호작용의 자연스러운 포착이라는
세 요구사항을 동시에 충족한다.
기존의 클리핑(clipping) 기반 접근이나 경계 미보장 모델과 달리,
시그모이드 바운딩은 점근적으로 경계에 도달하되 경계값 자체를 취하지 않아
수치적 안정성과 이론적 엄밀성을 동시에 확보한다.}


%% -------------------------------------------------------------------
\chg{둘째 성과는 \citet{lee2026drto}의 $\kappa^{*} = 1.2$에 대한 \frev{분포 환경 적합성} 종합 명제 도출이다.}
\label{subsec:achiev-kappa}
\chg{\citet{lee2026drto}의 $\kappa^{*}=1.2$를 본 연구의 새로운 분포 환경(C2 가우시안과 C3 파레토)에 적용한 결과, 9개 가설(H2--H10) PASS의 통합 해석을 통해 $\kappa = 1.2$의 분포 환경 \frev{운영 적합성}을 시사하였다(상세는 \secref{sec:ch4_hypothesis_summary}의 \frev{분포 환경 적합성} 종합 명제와 \subsecref{subsec:academic-methodological}).}

\fdel{시그모이드 곡률 파라미터 $\kappa$의 최적값을 다중 기준 복합 점수(multi-criteria
composite score) 방법론으로 도출하였다.
민감도(sensitivity), 선형성(linearity), 경계 여유(boundary margin),
해석 가능성(interpretability)의 4개 기준을 정규화하여 종합 평가한 결과,
$\kappa^{*} = 1.2$가 최적점으로 결정되었다.
$\kappa = 1.2$에서 관측성 채널의 $z_O \in [0,\; 1.2]$는
시그모이드의 비선형 압축이 의미 있게 작동하는 영역에 해당하여
($S(1.2) = 0.537$), 중간 입력에서의 감도를 확보하면서도
극단 영역에서의 비선형 포화(saturation) 거동이 물리적 경계 보장을
자연스럽게 실현한다.}


%% -------------------------------------------------------------------
\chg{셋째 성과는 C0$\to$C3 4단계 점진적 프레임워크와 \jrev{10개} 가설의 전원 채택이다.}
\label{subsec:achiev-framework}
\chg{C0--C3 4단계 점진적 조건과 6단계 다층 검증을 통해 \jrev{10개} 가설(H1--\jrev{H10})을 전원 채택하였다(상세는 \subsecref{subsec:academic-empirical}).}

\fdel{정적 기준선(C0), 관측성 도입(C1), 가우시안 불확실성(C2),
파레토 불확실성(C3)의 4단계 점진적 조건 체계를 수립하고,
6단계 다층 검증 프레임워크(L0 기반 검증(Golden Compare),
시뮬레이션 검증(H1--\jrev{H3}), 회귀 검증(\jrev{H4}--\jrev{H5}), 강건성 검증(\jrev{H6}--\jrev{H7}),
이중채널 검증(\jrev{H8}), 전문가 평가(\jrev{H9}--\jrev{H10}))를 통해 \jrev{10개} 가설을 체계적으로 검증하였다.
정량적 결과의 핵심은 다음과 같다.
C1에서 관측성에 의한 $\eta$ 감소($\bar{\eta} = 0.853$, Cohen's $d = -1.180$),
C2에서 불확실성 프리미엄의 발현($\bar{\eta} = 1.682$, Cohen's $d = +1.871$),
C3에서 평균 유사$\cdot$꼬리 급증($\bar{\eta} = 1.514$, CVaR$_{99}$ $+24.0\%$).
$R^2_{(2\mathrm{q})}$는 C1$\,= 1.000$, C2$\,= 0.998$이며,
C3는 전체 $R^2_{(2\mathrm{q})} = 0.858$이나 P85.8 분할 후 Core $R^2 = 0.999$로 \jrev{H4}가 지지되었다.}


%% -------------------------------------------------------------------
\chg{넷째 성과는 꼬리 감쇠 현상의 발견($0.52$배, 꼬리 영역)이다.}
\label{subsec:achiev-dual-channel}
\chg{P85.8 분할 회귀에서 $k_O$ 계수가 Core $-0.798$에서 Tail $-0.415$로 $0.52$배 감쇠하였고, CVaR$_{99}$는 $+24.0\%$ 증가하였다(상세는 \subsecref{subsec:academic-theoretical}).}

\fdel{C3 조건에서 P85.8 기준으로 분할 회귀를 수행한 결과,
관측성 가중치 $k_O$의 회귀 계수가 핵심 영역(Core, $-0.798$)에서
꼬리 영역(Tail, $-0.415$)으로 $0.52$배 감쇠되는 꼬리 감쇠
(tail attenuation) 현상을 발견하였다.
이 발견은 극단적 불확실성 환경에서 시그모이드 포화에 의해
관측성 투자의 한계 효용이 완화된다는 이론적 시사점을 제공한다.
P85.8 CDF 교차점은 가우시안(C2)과 파레토(C3)의 CDF가
교차하는 구조적 경계로서, 교차점 이하에서 두 분포는
사실상 구분 불가능하나(Cohen's $d \approx 0$),
교차점 초과에서 C3의 $\eta$가 C2를 급격히 상회하여
CVaR$_{99}$ 기준 $+24.0\%$의 격차를 보인다.}


%% -------------------------------------------------------------------
\chg{다섯째 성과는 전문가 검증(4.46/5.0)과 6단계 다층 검증이다.}
\label{subsec:achiev-validation}
\chg{전문가 5명의 평균이 4.46/5.0($p=0.011$, $d=1.65$)이고 Cronbach's $\alpha=0.89$, S-CVI/Ave $=0.985$였으며, 인터뷰에서 15개 주제 모두 $\geq 80\%$ 동의를 얻었다(상세는 \subsecref{subsec:academic-empirical}).}

\fdel{BCM/DR 분야 전문가 5명(평균 경력 15.2년)을 대상으로
수행한 정량적 설문에서 전체 평균 4.46/5.0점($t(4) = 3.68$, $p = 0.011$,
Cohen's $d = 1.65$)을 달성하였으며,
Cronbach's $\alpha = 0.89$, S-CVI/Ave$\,= 0.985$로
높은 내적 일관성과 평가자 간 일치도가 확인되었다(\jrev{H9} 지지).
심층 인터뷰에서는 5대 주제 15개 하위 주제 모두(100\%)에서 80\% 이상
전문가 동의가 달성되어 \jrev{H10}이 지지되었다.
시뮬레이션 검증, 회귀 검증, 강건성 검증, 이중채널 검증, 전문가 평가의
6단계(L0--L5)가 계층적 의존 구조를 형성하여 모델의 타당성을 다각적으로 보장하며,
정량적 시뮬레이션과 정성적 전문가 평가의 삼각검증(triangulation)을 통해
방법론적 편향을 최소화하였다.}

\vspace{0.5em}

\p{[}표~\vref{tab:ch7-findings-summary}\p{]}는 이상의 5대 성과를 핵심 결과,
정량적 근거, 학술적 의의의 관점에서 종합한 것이다.

\begin{table}[htbp]
\centering
\caption{핵심 연구 결과 요약}
\label{tab:ch7-findings-summary}
\small
\begin{tabularx}{\textwidth}{l X X}
\toprule
\textbf{핵심 결과} & \textbf{정량적 근거} & \textbf{학술적 의의} \\
\midrule
개별 시그모이드 바운딩 모델
  & $\DRTO \in (0,\; R_{\max})$ 보장, $C^{\infty}$ 연속
  & BCM 분야 최초의 물리적 경계 보장 동적 RTO 모델 \\
\addlinespace
\jrev{$\kappa^{*}=1.2$ \frev{분포 환경 적합성}}
  & \jrev{9개 가설(H2--H10) PASS 통합 해석}
  & \jrev{\chg{\citet{lee2026drto}의} 매개변수의 본 연구 환경 적용성을 다중 가설 PASS로 \frev{시사한} 학술적 명확성} \\
\addlinespace
C0$\to$C3 프레임워크, H1--\jrev{H10} 전원 채택
  & $\eta$: 0.853(C1), 1.682(C2), 1.514(C3); C1/C2 $R^2 \geq 0.998$, C3 Core $R^2 = 0.999$
  & 4단계 점진적 BCM 성숙도 모형의 실증적 근거 \\
\addlinespace
꼬리 감쇠 ($0.52$배)
  & Tail $k_O$ 계수 $-0.415$ vs Core $-0.798$; CVaR$_{99}$ $+24.0\%$
  & 극단 환경에서 관측성 투자 효용 감소의 최초 정량적 규명 \\
\addlinespace
전문가 검증 + 6단계 다층 검증
  & 4.46/5.0, $\alpha = 0.89$, S-CVI/Ave$\,= 0.985$; \jrev{10/10} 가설 채택
  & \chg{시뮬레이션과 전문가} 삼각검증의 방법론적 표준 제시 \\
\bottomrule
\end{tabularx}
\end{table}

\chg{이상의 5대 성과는 학술적 의의에 그치지 않고 본 장의 실무적 기여로 직접 이어진다. 개별 시그모이드
바운딩 모델과 회귀 설명력(첫째와 셋째 성과)은 회귀식 기반 $\DRTO$ 산출 도구(\subsecref{subsec:practical-regression-tool})로,
꼬리 감쇠 발견(넷째 성과)은 이중채널 감쇠 기반 자원배분과 스트레스 테스트
도구(\subsecref{subsec:practical-dualchannel-tool})로, $\kappa = 1.2$의 적합성과 전문가 검증(둘째와 다섯째 성과)은
산업별 적용 가이드라인과 운영 도구의 신뢰 근거로 활용된다. 한편 위 표가 다루는 다섯 핵심 성과 외에도
4장에서 도출된 세부 결과들이 의미의 크고 작음과 무관하게 각각 어느 학술 기여 또는 실무 활용으로
귀착되는지는 \subsecref{subsec:ch5_results_map}의 전수 매핑에 정리되어 있다.}


